\chapter{IMPLEMENTATION AND EXPERIMENTAL RESULTS}
\label{chap:implementation}

[This chapter details the development environment, implementation highlights, experimental setup, and analysis of results.]

% ================= SECTION 4.1 =================
\section{Development Environment}
\label{sec:environment}

[Describe the hardware and software configuration used during development.]

\begin{table}[H]
    \centering
    \caption{Experimental Hardware and Software Environment}
    \label{tab:sample_env}
    \begin{tabular}{|l|l|}
        \hline
        \textbf{Component} & \textbf{Specification} \\ \hline
        Operating System & Linux (Ubuntu / Arch Linux) \\ \hline
        Programming Language & Python 3.x \\ \hline
        Compiler / Toolchain & GCC / Clang / Python runtime \\ \hline
        Memory (RAM) & 16 GB \\ \hline
    \end{tabular}
    \source{Author's testbed setup}
\end{table}

% ================= SECTION 4.2 =================
\section{Implementation Details}
\label{sec:code_details}

A sample algorithm implementation in Python is shown in Listing~\ref{lst:sample_code}.

\begin{lstlisting}[language=Python, caption={Sample Algorithm Implementation}, label={lst:sample_code}]
def compute_metrics(data_samples, threshold=0.5):
    """
    Compute classification evaluation metrics.
    """
    results = []
    for sample in data_samples:
        score = sample.get("score", 0.0)
        is_positive = score >= threshold
        results.append({"id": sample["id"], "prediction": is_positive})
    return results

if __name__ == "__main__":
    test_data = [{"id": 1, "score": 0.82}, {"id": 2, "score": 0.35}]
    print(compute_metrics(test_data))
\end{lstlisting}

% ================= SECTION 4.3 =================
\section{Experimental Results and Discussion}
\label{sec:results}

[Present quantitative results, comparative performance, and analytical discussion.]

\begin{table}[H]
    \centering
    \caption{Experimental Performance Evaluation}
    \label{tab:sample_results}
    \begin{tabular}{|c|c|c|c|}
        \hline
        \textbf{Test Case} & \textbf{Data Size} & \textbf{Execution Time (ms)} & \textbf{Status} \\ \hline
        TC-01 & 1 KB & 0.45 & Passed \\ \hline
        TC-02 & 10 KB & 2.18 & Passed \\ \hline
        TC-03 & 100 KB & 18.64 & Passed \\ \hline
    \end{tabular}
    \source{Experimental measurements recorded by the author}
\end{table}
